\documentclass[11pt]{article}
\usepackage[margin=0.86in]{geometry}
\usepackage[T1]{fontenc}
\usepackage{lmodern}
\usepackage{amsmath,amssymb,amsthm,graphicx,microtype}
\usepackage[hidelinks]{hyperref}

\newtheorem{theorem}{Theorem}
\newtheorem{lemma}[theorem]{Lemma}
\newcommand{\R}{\mathbb R}
\newcommand{\N}{\mathbb N}
\newcommand{\cB}{\mathcal B}
\newcommand{\cX}{\mathcal X}
\newcommand{\fc}{\mathfrak c}

\title{A Maximal-Disjoint-Family Argument Refutes\\
the Infinite-Lebesgue-Measure Independence Conjecture}
\author{Candidate counterexample packet\\
\small arXiv:2503.06217, Conjecture 4.6}
\date{11 August 2026}

\begin{document}
\maketitle

\noindent\textbf{Status.} Candidate full ZFC disproof, subject to expert
review.  In fact the argument constructs the stronger object that the source
calls a $\fc$-separable \emph{partition}; hence the conjectured failure of a
$\fc$-separable envelope is impossible.

\section{Source question and exact consequence}

Rodríguez-Vidanes and Sampedro study Baker's infinite-dimensional Lebesgue
measure $\mu$ on the Borel space $(\R^{\N},\cB_\infty)$
\cite{source}.  Their Conjecture 4.6, printed on page 19, reads:
\begin{quote}
There is a model of ZFC in which $(\R^{\N},\mu)$ does not admit a
$\fc$-separable envelope.
\end{quote}

\begin{center}
\includegraphics[width=\textwidth]{figures/open_problem_crop.png}
\end{center}

The authors state that this would make
\[
 L^p(\mu)\cong \ell^p(\fc,L^p[0,1])\qquad(p\ne2)
 \tag{1}
\]
independent of ZFC.  We prove instead that ZFC itself gives the requisite
partition and therefore proves (1), for every $1\le p<\infty$.

\section{Definitions and source inputs}

In the terminology of \cite[Definition 4.2]{source}, a family
$\{X_i:i\in I\}$ is a $\kappa$-separable partition if $|I|=\kappa$, every
$X_i$ is a nonzero, sigma-finite, separable measurable set, distinct members
intersect in a null set, and
\[
 \nu(E)=\sum_{i\in I}\nu(E\cap X_i)                         \tag{2}
\]
for every measurable $E$ of finite measure.  The sum over an arbitrary index
set means the supremum of its finite subsums.  A separable envelope only asks
for (2) when $E$ is contained in the union of the $X_i$; thus a separable
partition is stronger.

We use three facts proved in the source:
\begin{enumerate}
\item $\cB_\infty$ is generated by countably many rational cylinder sets, and
every finite-measure measurable set is separable
\cite[Lemma 2.3]{source};
\item Section 3 constructs $\fc$ many measurable rectangles $C_A$, each of
measure one, with $\mu(C_A\cap C_B)=0$ for $A\ne B$
\cite[proof of the main theorem of Section 3]{source};
\item $(\R^{\N},\mu)$ is relatively nonatomic
\cite[Appendix A]{source}, and a relatively nonatomic space with a
$\kappa$-separable partition has
$L^p(\nu)\cong\ell^p(\kappa,L^p[0,1])$
\cite[Theorem 4.3]{source}.
\end{enumerate}

\section{Proof intuition}

The only apparent difficulty is that an uncountable almost-disjoint family
need not have a measurable union.  Maximality avoids taking that union.  Fix
a finite-measure test set $E$.  Among a pairwise-null family, only countably
many members can capture positive measure from $E$.  Remove precisely those
countably many pieces.  If a positive remainder survived, it would itself be
a finite-measure set almost disjoint from every member of the maximal family,
contradicting maximality.  This proves the full partition identity locally on
each finite $E$.

The source's continuum-sized rectangle family supplies the lower cardinal
bound.  Countable generation of the ambient sigma-algebra supplies the upper
bound.  No form of CH is used.

\section{The maximal-disjointization lemma}

\begin{lemma}\label{lem:maximal}
Let $(X,\Sigma,\nu)$ be a measure space, and let $\cX$ be a maximal family of
measurable sets such that
\[
 0<\nu(X_i)<\infty\quad(X_i\in\cX),
 \qquad \nu(X_i\cap X_j)=0\quad(X_i\ne X_j).
\]
Then for every $E\in\Sigma$ with $\nu(E)<\infty$,
\[
 \nu(E)=\sum_{X_i\in\cX}\nu(E\cap X_i).                   \tag{3}
\]
\end{lemma}

\begin{proof}
For every finite $F\subset\cX$, pairwise nullity gives
\[
 \sum_{X_i\in F}\nu(E\cap X_i)\le \nu(E).
\]
Let
\[
 J_E:=\{X_i\in\cX:\nu(E\cap X_i)>0\}.
\]
This set is countable.  Indeed, for every $n\ge1$ there can be only finitely
many $i$ with $\nu(E\cap X_i)\ge1/n$; otherwise finite subsums would exceed
$\nu(E)$.  Every positive number is at least $1/n$ for some $n$, so $J_E$ is
a countable union of finite sets.

Enumerate $J_E$ (with the finite case included) and use countable additivity
after discarding the pairwise-null overlaps.  Then
\[
 \nu\!\left(E\cap\bigcup_{X_i\in J_E}X_i\right)
   =\sum_{X_i\in J_E}\nu(E\cap X_i).
\]
If the right-hand side were strictly smaller than $\nu(E)$, the measurable
set
\[
 F_E:=E\setminus\bigcup_{X_i\in J_E}X_i
\]
would satisfy $0<\nu(F_E)<\infty$.  It is disjoint from each member of $J_E$.
For $X_i\notin J_E$, we have
$\nu(F_E\cap X_i)\le\nu(E\cap X_i)=0$.  Thus $F_E$ could be adjoined to
$\cX$, contradicting maximality.  Hence equality holds, which is (3), since
all terms outside $J_E$ vanish.
\end{proof}

\section{Main result}

\begin{theorem}\label{thm:main}
ZFC proves that $(\R^{\N},\cB_\infty,\mu)$ admits a
$\fc$-separable partition.  Consequently Conjecture 4.6 of \cite{source} is
false, and
\[
 L^p(\mu)\cong\ell^p(\fc,L^p[0,1])
\]
isometrically for every $1\le p<\infty$, without CH.
\end{theorem}

\begin{proof}
Let $\mathcal A_0=\{C_A:A\in\mathcal A\}$ be the source's continuum-sized family of
measure-one rectangles with pairwise null intersections.  Order by inclusion
the families of finite positive measurable sets which contain $\mathcal A_0$ and
whose distinct members have null intersection.  The union of every chain is
again such a family, so Zorn's lemma gives a maximal extension $\cX$.

Every member of $\cX$ has finite positive measure, hence is sigma-finite; by
source Lemma 2.3 it is separable.  Lemma~\ref{lem:maximal} gives identity (2)
for every finite-measure $E$.  Thus $\cX$ is a separable partition apart from
checking its cardinal.

We have $|\cX|\ge|\mathcal A_0|=\fc$.  On the other hand, the countably generated
sigma-algebra $\cB_\infty$ has cardinal at most
$2^{\aleph_0}=\fc$, so $|\cX|\le\fc$.  Hence $|\cX|=\fc$ in ZFC.  This proves
the first assertion.  The isometry follows immediately from relative
nonatomicity and source Theorem 4.3.
\end{proof}

\section{Verification, novelty, and review focus}

The proof was audited in five independent places: existence of the maximal
extension; countability of $J_E$; measurability and measure of the remainder;
admissibility of each maximal-family member; and the exact cardinal
$|\cX|=\fc$.  Notice especially that no uncountable union is asserted to be
measurable.  The argument proves the source's stronger partition definition,
so no ambiguity in its weaker envelope definition affects the result.

A bounded novelty search through 11 August 2026 covered the run's solution,
attempt, proof-gap, and registry indexes; the exact wording of Conjecture 4.6;
the title and arXiv id; the DOI; author names; and combinations of
``separable envelope,'' ``maximal disjoint family,'' and ``ZFC.''  The current
arXiv PDF and published metadata were also checked.  No correction, erratum,
or later resolution was found.  This is evidence about retrieval, not proof
of novelty.

\medskip
\noindent\textbf{Human-review recommendation.} High priority.  The conclusion
is unexpectedly stronger than the source's CH theorem, while the proof is
short.  Review should focus on whether ``separable subset'' in Definitions
4.2 and 4.4 is exactly the separability property supplied by Lemma 2.3, and
on the remainder step in Lemma~\ref{lem:maximal}.  Under the source's printed
definitions, both checks are direct.

\begin{thebibliography}{9}
\bibitem{source}
D. L. Rodríguez-Vidanes and J. C. Sampedro,
\emph{Isometric classification of the $L^p$-spaces of infinite dimensional
Lebesgue measure}, arXiv:2503.06217; Banach J. Math. Anal. \textbf{20}
(2026), article 7, DOI: \href{https://doi.org/10.1007/s43037-025-00473-y}
{10.1007/s43037-025-00473-y}.
\end{thebibliography}

\end{document}
